% !TeX root = ../main.tex

\chapter{Code and Data Availability}
\label{app:availability}

\section{Code Availability}
\label{sec:code-availability}

All source code, model definitions, training and evaluation scripts, and the complete set of experiment configuration files are publicly available at:

\begin{center}
\url{https://github.com/m2lab-ntu/gfm-classifier}
\end{center}

The repository contains, in addition to the training and inference pipeline, the full YAML configuration for every experiment reported in Chapter~\ref{ch:experiments} (the \texttt{v8}--\texttt{v15} runs spanning the 5M--250M read scales, the from-scratch and warm-start NT-v2 variants, and the MetaTransformer 13-mer baseline), under \texttt{configs/}. Each configuration records the backbone, the attention-pooling head dimensions, the LoRA settings ($r{=}16$, $\alpha{=}32$, applied to the query/key/value projections), the optimisation schedule, and the reverse-complement test-time augmentation flag, so that any reported result can be reproduced from the corresponding config and checkpoint. The headline hyperparameters are summarised in Tables~\ref{tab:lora-config} and~\ref{tab:training-config} of Chapter~\ref{ch:methodology}.

The repository also provides the distributed-training launch scripts used for the large-scale (250M-read) runs, which employ multi-GPU \texttt{DistributedDataParallel} rather than the single-device setting used for the smaller experiments.

\section{Data Availability}
\label{sec:data-availability}

The reference genome set (1{,}535 microbial genomes) is drawn from publicly available assemblies; the accession list, the read-simulation procedure, and the scripts used to construct the balanced and natural-abundance read sets (including the unique-read accounting for the 50M and 250M training pools) are included in the repository above under \texttt{data/}. Processed read sets and trained checkpoints are available from the author on reasonable request.
